\chapter{Hereditary Non-Polyposis Colorectal Cancer (HNPCC)}

\section*{Definition and Overview}
Hereditary non-polyposis colorectal cancer (HNPCC), now more commonly called Lynch syndrome, is an inherited condition that greatly increases the risk of colorectal cancer and several other cancers, including cancer of the uterus, ovaries, stomach, and urinary tract. It is caused by changes (mutations) in genes that normally repair mistakes in DNA, such as MLH1, MSH2, MSH6, and PMS2. People with Lynch syndrome have around a 40--70\% lifetime risk of colorectal cancer, often at a younger age than the general population. Surveillance and risk-reducing surgery dramatically improve outcomes.

\section*{Epidemiology}
Lynch syndrome affects about one in 280 to one in 400 people, making it the most common inherited colorectal cancer syndrome. It accounts for around 3--5\% of all colorectal cancers and a higher proportion of those diagnosed under 50. An estimated 1 in 10 Australians with a strong family history of bowel cancer carries the gene change. Because the risk of cancer is high and onset can be early, identifying families through genetic testing and surveillance saves lives.

\section*{Aetiology and Pathophysiology}
\begin{itemize}
    \item \textbf{Genetic cause:} inherited mutations in DNA mismatch repair genes (MLH1, MSH2, MSH6, PMS2, and EPCAM)
    \item \textbf{Inheritance:} autosomal dominant, so each child of an affected parent has a 50\% chance of inheriting the mutation
    \item \textbf{Pathophysiology:} the mismatch repair system normally fixes errors made during DNA replication; when it fails, mutations accumulate rapidly, particularly in sections of DNA with repeated sequences (microsatellite instability)
    \item \textbf{Cancer risks:} colorectal cancer 40--70\%, endometrial cancer 40--60\% in women, plus ovarian, stomach, small bowel, urinary tract, pancreatic, and other cancers
    \item \textbf{Early onset:} bowel cancer often develops before age 50
    \item \textbf{Tumour testing:} bowel tumours from Lynch syndrome typically show microsatellite instability and loss of mismatch repair proteins
\end{itemize}

\section*{Clinical Features}
\begin{itemize}
    \item A personal or family history of colorectal or endometrial cancer at a young age
    \item Multiple family members on the same side of the family with Lynch syndrome cancers
    \item A personal history of multiple bowel cancers or bowel cancer with another Lynch cancer
    \item The condition itself causes no symptoms; features emerge through cancer
    \item Blood in the stool, bowel habit change, or anaemia may indicate bowel cancer
    \item Abnormal uterine bleeding may indicate endometrial cancer
\end{itemize}

\section*{Differential Diagnosis}
\begin{itemize}
    \item Familial adenomatous polyposis (FAP), a different inherited syndrome with hundreds of polyps
    \item Other hereditary syndromes, including MUTYH-associated polyposis and hamartomatous polyposis
    \item Sporadic bowel cancer in a family, which is common without a hereditary syndrome
    \item Constitutional mismatch repair deficiency, a rare childhood syndrome
    \item Genetic testing and tumour testing distinguish Lynch syndrome from these conditions
\end{itemize}

\section*{When to Seek Medical Care}
Anyone with a personal or family history suggesting Lynch syndrome should discuss referral to a familial cancer service. National criteria and predictive tools identify who should be offered genetic testing.

\textbf{Call 000 (or local emergency services) for:}
\begin{itemize}
    \item Heavy rectal bleeding with faintness or dizziness
    \item Sudden severe abdominal pain, which may indicate obstruction or perforation
    \item Vomiting blood or passing black stools
    \item Severe pelvic pain or heavy vaginal bleeding
\end{itemize}

\section*{Diagnosis and Investigations}
\begin{itemize}
    \item \textbf{Tumour testing:} colorectal tumours are screened for microsatellite instability and mismatch repair protein loss, with every new bowel cancer ideally tested
    \item \textbf{Genetic testing:} blood tests identify mutations in the mismatch repair genes, confirming the diagnosis
    \item \textbf{Predictive testing:} family members can be tested for the specific family mutation
    \item \textbf{Genetic counselling:} essential before and after testing to discuss implications
    \item \textbf{Family history assessment:} structured tools identify people at high risk
\end{itemize}

\section*{Management}
\begin{itemize}
    \item \textbf{Bowel surveillance:} colonoscopy every 1--2 years from age 25, removing polyps before they become cancerous
    \item \textbf{Endometrial surveillance:} for women, including gynaecological examination and ultrasound, with endometrial sampling
    \item \textbf{Other surveillance:} upper gastrointestinal endoscopy and urine testing at intervals
    \item \textbf{Preventive surgery:} colectomy and hysterectomy with removal of the ovaries are options for people at highest risk
    \item \textbf{Aspirin:} may reduce the risk of bowel cancer in Lynch syndrome and is discussed with specialists
    \item \textbf{Family communication:} informing relatives so they can be tested and screened
    \item \textbf{Reproductive options:} pre-implantation genetic diagnosis for couples planning children
\end{itemize}

\section*{Complications}
\begin{itemize}
    \item Colorectal cancer, often before the age of screening in the general population
    \item Endometrial, ovarian, and other Lynch-associated cancers
    \item Metachronous cancers: people who survive one cancer have an elevated risk of second cancers
    \item Psychological impact of genetic risk and screening burden
    \item Surgical complications of preventive surgery
    \item Non-adherence to surveillance, which removes the protective benefit
\end{itemize}

\section*{Prognosis and Outcomes}
With regular surveillance, the outlook for people with Lynch syndrome is excellent: colonoscopic surveillance reduces colorectal cancer risk and deaths by more than half, and most bowel cancers in screened individuals are found early and cured. Women who choose risk-reducing gynaecological surgery virtually eliminate their endometrial and ovarian cancer risk. Genetic testing also allows unaffected family members to be reassured or managed appropriately. Lynch syndrome is one of the best examples of how genetic information translates directly into life-saving prevention.

\section*{Prevention and Self-Care}
\begin{itemize}
    \item Attend all recommended surveillance tests on time
    \item Discuss genetic testing with a familial cancer service if you have a family history
    \item Share genetic information with close relatives so they can be tested
    \item Report new bowel or gynaecological symptoms promptly
    \item Maintain a healthy lifestyle: fibre-rich diet, activity, healthy weight, and limited alcohol
    \item Discuss aspirin prevention with your specialist
    \item Consider preventive surgery after informed discussion if your risk is high
    \item Seek psychological support to manage the anxiety of genetic risk
\end{itemize}

\section*{Sources and Further Reading}
The following reputable sources provide further information for healthcare students and patients seeking reliable, current advice about Lynch syndrome (HNPCC).

\begin{itemize}
    \item Lynch Syndrome Australia. \textit{Information and support for families.} https://lynchsyndromeaustralia.org.au/
    \item Bowel Cancer Australia. \textit{Hereditary bowel cancer and Lynch syndrome.} https://www.bowelcanceraustralia.org/
    \item eviQ Cancer Genetics. \textit{Lynch syndrome risk management.} https://www.eviq.org.au/
    \item Cancer Australia. \textit{Genetic testing for hereditary cancer.} https://www.canceraustralia.gov.au/
    \item NHS. \textit{Lynch syndrome: overview and screening.} https://www.nhs.uk/conditions/lynch-syndrome/
    \item National Cancer Institute. \textit{Lynch syndrome genetics.} https://www.cancer.gov/
\end{itemize}
